\documentclass{article} % amsmart?

\usepackage{amsmath}
\usepackage{amsthm}
\usepackage{graphicx}
\usepackage{amssymb}

\newcommand{\emptyline}{\vspace{\baselineskip}}

\title{Problem 5}
\author{Caleb Hessing}
\date{Oct 7, 2026}

\begin{document}

\maketitle

\begin{center}
\textit A: For all real numbers $x$, there exists an integer $n$ such that $n \le x < n+1$.
\textit B: There exists an integer $n$ such that, for all real numbers $x$, $n \le x < n+1$.
\end{center}

\textbf{Part (a)} 

\[\begin{gathered} 
  \textit{A: } \forall x \in \mathbb R, \exists n \in \mathbb Z : n \le x < n + 1. \\
  \textit{B: } \exists n \in \mathbb Z : n \le x < n + 1 \nobreakspace \forall x \in \mathbb R.  
\end{gathered}\]

\emptyline
\textbf{Part (b)} 

\textit A is saying that all real numbers are bounded between some two consecutive integers. This statement is true.
\textit B is saying that there are consecutive integers between which all real numbers lie. This statement is not true.

\emptyline
\textbf{Part (c)}

Suppose every person has exactly one soulmate, and that everyone is the soulmate of their soulmate. 
Consider then these two statements:

\begin{center}
    \textit {A: } For every person $x$, there exists another person $y$ \\ such that $x$ and $y$ are soulmates. \\
    \textit {B: } There exists some person $y$ such that, for every person $x$, \\ $x$ and $y$ are soulmates.
\end{center}

Statement A is true. It reflects the (unequivocally) true property that everyone has a soulmate.
However, statement B suggests that there is some person that is \textit everyone's soulmate.
But every person (undoubtedly) only has one soulmate, so this is false. 

\end{document}